\chapter{Evaluating a Hybrid Planner: Experiments, Metrics and Benchmarks}
\label{ch:ch25}
% Function names for pseudocode are prefixed with "Eval" to stay unique
% across the book.
\SetKwFunction{EvalPairedStudy}{PairedStudy}
\SetKwFunction{EvalScenario}{MakeScenario}
\SetKwFunction{EvalRun}{RunEpisode}
\SetKwFunction{EvalMetrics}{Metrics}
\SetKwFunction{EvalCompare}{PairedCompare}
\SetKwFunction{EvalBootstrap}{BootstrapInterval}

\begin{objectives}
\begin{itemize}
  \item Define the ten metrics of the training plan precisely, with a formula, a unit, an estimator over runs and a list of what each one is sensitive to, and sort them into safety, efficiency, effort and constraint metrics.
  \item Turn the five factors of the plan into an experiment matrix, choose between a full factorial and a fractional design, and name the baselines that every comparison needs.
  \item Generate controlled scenarios from seeds, with intruder trajectory families and difficulty knobs, so that every strategy sees exactly the same scenarios.
  \item Choose the right statistics: how many runs, mean or median, confidence intervals, paired comparisons, effect sizes, success rates under a cap, and the three plots that show them.
  \item Run a small paired study with the book's code, read its table and figures cautiously, and write a results section that a reviewer can reproduce and trust.
\end{itemize}
\end{objectives}

% ---------------------------------------------------------------------
\section{Why this matters for a drone swarm}
\label{sec:ch25-motivation}

Every algorithm in this book is established. \astar, \dstarlite, \cbs and
\ecbs, \orca, DWA, the Kalman filter and MPC each have their original paper,
their proofs and their decades of use, and the training plan says so plainly:
the contribution of your research will not be a new planner built from zero.
It will be a \emph{combination}: the hybrid architecture of \cref{ch:ch24},
which lets a conflict-free global plan, a predictor for an unknown drone, a
reactive layer and a replanner share one vehicle. A combination has no
theorem that says it is better than its parts. The only way to know whether
the hybrid planner is safer than local avoidance alone, or cheaper than
replanning everything, or whether a learned predictor earns its cost, is to
measure it --- and to measure it in a way that another researcher can repeat
and believe.

This is where research contributions are made or lost. A reviewer who reads
that ``the hybrid planner reduced collisions by 40\,\%'' will ask: on which
scenarios, against which baselines, over how many runs, with what spread, and
did the baselines see the same scenarios? If the answers are missing, the
claim is worth nothing, however good the code. Conversely, a modest result
reported with its scenarios, seeds, intervals and failure cases is a result
that can be built on. The last row of the training plan, Week~12
(\cref{ch:appA}), asks for exactly this: ``a research prototype and an
experiment matrix suitable for a paper or thesis chapter''.

The chapter proceeds as follows. \Cref{sec:ch25-intuition} gives the shape of
an evaluation pipeline in one picture. \Cref{sec:ch25-metrics} defines the
metrics, \cref{sec:ch25-design} the experiment matrix and the baselines,
\cref{sec:ch25-scenarios} the scenario generator and the paired design that
makes comparisons fair. \Cref{sec:ch25-statistics} covers the statistics and
the plots, each shown on data produced by the book's own toy simulator.
\Cref{sec:ch25-study} runs a complete mini-study (two and four drones, zero or
one intruder, three strategies, twenty seeds) and interprets it.
\Cref{sec:ch25-reproducibility,sec:ch25-reporting} say what to save and what
to write, \cref{sec:ch25-benchmarks} surveys the benchmarks and physics
simulators you can connect to, and \cref{sec:ch25-implementation} discusses
the code. Unlike the algorithm chapters, this one has no theorem about a
planner; its two small propositions are about experiments.

% ---------------------------------------------------------------------
\section{The idea in plain words}
\label{sec:ch25-intuition}

An evaluation is a pipeline (\cref{fig:ch25-pipeline}). A
\textbf{scenario generator}\index{scenario generator} turns a seed and a
cell of the experiment matrix into a concrete world: map, starts, goals,
intruder trajectories, even the measurement noise the predictor will see. A
\textbf{runner} executes every strategy on that same scenario and writes a
log of positions, decisions and timings. \textbf{Metric functions} reduce
each log to a row of numbers. \textbf{Statistics} reduce the rows of a cell
to estimates with intervals and compare strategies \emph{pairwise, scenario
by scenario}. Only then comes the \textbf{report}: tables, three kinds of
plots, and the limitations. Nothing in the pipeline is clever; what makes it
scientific is that every arrow is deterministic given the seed, so that the
whole picture can be regenerated by anyone who has the code and the
configuration file.

\begin{figure}[tb]
  \centering
  \inputfigure{ch25/pipeline}
  \caption{The evaluation pipeline. The seed and the cell of the experiment
  matrix fix the scenario; every strategy runs on the same scenario (paired
  design); logs become metric rows, rows become estimates with intervals, and
  only the last box is prose. The dashed arrow is the one to avoid: tuning a
  strategy on the scenarios that will be reported.}
  \label{fig:ch25-pipeline}
\end{figure}

\begin{keyidea}
Fix the scenarios by seeds, run every strategy on the same scenarios, report
every metric with its spread and its failures, and let the seeds regenerate
the whole study. A difference you cannot show on paired runs with a
confidence interval is not a result yet.
\end{keyidea}

% ---------------------------------------------------------------------
\section{Metrics: what to measure and how}
\label{sec:ch25-metrics}

\subsection{The execution log}

All metrics are functions of one object, the log of a run.

\begin{definition}[Execution log]
\label{def:ch25-log}
A \textbf{run}\index{run (evaluation)} executes one scenario under one
strategy for $T+1$ sampling instants $t = 0, 1, \dots, T$ of length $\dt$. Its
\textbf{execution log}\index{execution log} records, for every $t$, the
positions $\pos_i(t) \in \R^2$ (or $\R^3$) of the $k$ controlled drones
$i = 1, \dots, k$, the positions $\vect{q}_j(t)$ of the $m$ intruders
$j = 1, \dots, m$, the state of each drone's state machine (Nominal, Avoiding,
Reconnecting or Replanning, \cref{ch:ch24}), the computation time $c_i(t)$
of the decision taken for drone $i$ at $t$, and the events of the run, in
particular every replan. It also stores the nominal plan
$\pos_i^{\mathrm{nom}}(t)$ of every drone, the goals $\vect{g}_i$ and the
radii $r_i$, so that the metrics can be recomputed later from the log alone.
\end{definition}

The separation of two agents at time $t$ is $d_{ij}(t) = \norm{\pos_i(t) -
\pos_j(t)}$ as in \cref{ch:ch02}, and two agents of radii $r_i$, $r_j$
collide when $d_{ij}(t) < R = r_i + r_j$. Positions are sampled, so a
collision between two samples can be missed. The code of this chapter
therefore evaluates the minimum of $d_{ij}$ on every interval $[t, t+1]$
under the assumption of linear motion, using the closest-approach formula of
\cref{ch:ch02}: with $\vect{r}_0 = \pos_i(t) - \pos_j(t)$ and
$\vect{w} = \pos_i(t+1) - \pos_j(t+1) - \vect{r}_0$, the minimum over the
interval is $\norm{\vect{r}_0 + s^*\vect{w}}$ with
$s^* = \operatorname{clip}(-\vect{r}_0 \cdot \vect{w} / \norm{\vect{w}}^2,
0, 1)$. With $\dt = 0.1$~s and relative speeds up to $4$~m/s a sampled
minimum can be off by $0.2$~m, a third of the collision distance of two
$0.3$~m drones.

\subsection{Four families of metrics}

The plan lists ten metrics. They answer four different questions, and a
results table should group them that way: did anything go wrong
(\textbf{safety}), how much did avoidance cost the mission
(\textbf{efficiency}), how hard did the planner work (\textbf{effort}) and
were the swarm constraints respected (\textbf{constraints}).

\begin{definition}[Safety metrics]
\label{def:ch25-safety}
The \textbf{collision indicator}\index{collision rate} of a run is
$X_{\mathrm{coll}} = 1$ if some pair of agents (drone--drone or
drone--intruder) has $d_{ij}(t) < r_i + r_j$ at some time, and $0$
otherwise; the \textbf{collision count} $n_{\mathrm{coll}}$ is the number
of such pairs. The \textbf{collision rate} of a cell is the fraction of its
$N$ runs with $X_{\mathrm{coll}} = 1$,
\begin{equation}
  \hat p_{\mathrm{coll}} = \frac{1}{N}\sum_{s=1}^{N} X_{\mathrm{coll}}^{(s)}.
  \label{eq:ch25-collrate}
\end{equation}
The \textbf{minimum separation}\index{minimum separation} of a run is
\begin{equation}
  d_{\min} = \min_{t}\ \min_{i \ne j} d_{ij}(t),
  \label{eq:ch25-dmin}
\end{equation}
reported separately for drone--drone pairs ($d_{\min}^{\mathrm{dd}}$) and
drone--intruder pairs ($d_{\min}^{\mathrm{di}}$), in metres. A run is a
\textbf{near miss}\index{near miss} if $d_{\min} < r_{\mathrm{safe}}$, the
safety radius the avoidance layers of \cref{ch:ch24} try to keep.
\end{definition}

\begin{definition}[Efficiency metrics]
\label{def:ch25-efficiency}
The \textbf{path length}\index{path length (metric)} of drone $i$ is the
flown length $L_i = \sum_{t=0}^{T-1} \norm{\pos_i(t+1) - \pos_i(t)}$ in
metres, and its \textbf{path-length ratio} is $\rho_i = L_i /
L_i^{\mathrm{nom}}$, where $L_i^{\mathrm{nom}}$ is the length of the nominal
plan; a run reports the mean ratio $\rho_L = \frac{1}{k}\sum_i \rho_i$. The
\textbf{travel time}\index{travel time} of drone $i$ is the first instant
after which it stays within the goal tolerance $\eps_g$ of $\vect{g}_i$,
\begin{equation}
  T_i = \dt \cdot \min\set{\tau : \norm{\pos_i(t) - \vect{g}_i} \le \eps_g
  \text{ for all } t \ge \tau},
  \label{eq:ch25-travel}
\end{equation}
in seconds, and $T_i = \infty$ if the drone is not at its goal at the end of
the log. The \textbf{makespan}\index{makespan!as a metric} and the
\textbf{sum of costs}\index{sum of costs!as a metric} of the run are
\begin{equation}
  \makespan = \max_{i} T_i, \qquad \sumcost = \sum_{i=1}^{k} T_i,
  \label{eq:ch25-objectives}
\end{equation}
the continuous-time forms of \cref{ch:ch07}: a MAPF path cost of $c$ steps
is a travel time of $c\,\dt$ seconds, and trailing waits at the goal are free
in both.
\end{definition}

\begin{definition}[Effort metrics]
\label{def:ch25-effort}
The \textbf{replanning count}\index{replanning count} $n_{\mathrm{rp}}$ is
the number of entries into the Replanning state summed over the drones of a
run; the number of steps spent in Avoiding or Reconnecting,
$n_{\mathrm{av}}$, is reported next to it. The \textbf{computation
time}\index{computation time} is recorded per decision, $c_i(t)$ in
milliseconds, and summarised by its mean $\bar c$, its 99th percentile
$c_{99}$ (the deadline behaviour), its maximum, the number of
\textbf{deadline misses} $\abs{\set{(i,t) : c_i(t) > \dt}}$, and the total
$C_{\mathrm{tot}} = \sum_{i,t} c_i(t)$ per run in seconds.
\end{definition}

\begin{definition}[Constraint metrics]
\label{def:ch25-constraints}
The \textbf{formation error}\index{formation error!as a metric} at time $t$
is $e_F(t)$ of \cref{ch:ch23}: the root-mean-square violation of the desired
displacements $\vect{d}_{ij}(t)$ over the edges of the formation graph $G_F$,
\begin{equation}
  e_F(t) = \sqrt{\frac{1}{\abs{E_F}}\sum_{\set{i,j} \in E_F}
    \norm{\pos_j(t) - \pos_i(t) - \vect{d}_{ij}(t)}^2},
  \label{eq:ch25-ef}
\end{equation}
where the template may be a fixed formation or, for a coordinated mission,
the relative geometry of the nominal plan, $\vect{d}_{ij}(t) =
\pos_j^{\mathrm{nom}}(t) - \pos_i^{\mathrm{nom}}(t)$. A run reports the
time average $\bar e_F$ and the peak $e_F^{\max}$, in metres. The
\textbf{communication violations}\index{communication violation!as a metric}
of a run are the number of time steps at which a kept edge is longer than
the communication range or the radius graph of \cref{ch:ch23} is
disconnected,
\begin{equation}
  n_{\mathrm{cv}} = \Bigl|\Bigl\{ t : \exists \set{i,j} \in E_{\mathrm{keep}}
  \text{ with } d_{ij}(t) > R_{\mathrm{comm}}, \text{ or }
  \lambda_2\bigl(\mat{L}(t)\bigr) = 0 \Bigr\}\Bigr|,
  \label{eq:ch25-commviol}
\end{equation}
where $\mat{L}(t)$ is the Laplacian of the radius graph at time $t$ and
$\lambda_2$ its algebraic connectivity.
\end{definition}

\subsection{Estimators, units and sensitivities}

A metric is a number per run; a result is a number per \emph{cell} of the
experiment matrix, estimated from the $N$ runs of the cell, and the estimator
must fit the metric. Rates such as \cref{eq:ch25-collrate} are proportions,
whose uncertainty is a binomial interval (the Wilson interval of
\cref{sec:ch25-statistics}). Lengths, times and errors are averaged with a
$t$ interval when their distribution is roughly symmetric, and reported as a
median with the interquartile range (IQR) when it is skewed. The minimum
separation is the worst case of a run, so its distribution across runs has a
heavy left tail; report its median, its 5th percentile and the near-miss
rate, never only its mean. Computation times are skewed to the right and
should be reported with their 99th percentile. \Cref{tab:ch25-metrics}
collects the ten metrics with their formulas, units, estimators and the
things they are sensitive to; the last column is the one to reread before
you believe a number.

\begin{table}[tb]
  \centering\footnotesize
  \caption{The ten metrics of the training plan. ``Estimator'' is the
  quantity reported for a cell of $N$ runs; ``sensitive to'' lists what can
  change the number without changing the planner.}
  \label{tab:ch25-metrics}
  \begin{tabularx}{\textwidth}{@{}L{2.4cm}L{3.2cm}L{1.1cm}L{2.9cm}Y@{}}
  \toprule
  Metric & Formula per run & Unit & Estimator over runs & Sensitive to\\
  \midrule
  \multicolumn{5}{@{}l}{\emph{Safety}}\\
  Collision rate & $X_{\mathrm{coll}} = [\exists\, t, i \ne j:\ d_{ij}(t) < r_i + r_j]$ & -- & proportion \cref{eq:ch25-collrate} with Wilson interval; $n_{\mathrm{coll}}$ alongside & sampling interval (missed contacts), radii, $N$ (granularity $1/N$), which pairs count\\
  Minimum separation & $d_{\min} = \min_{t,\,i\ne j} d_{ij}(t)$, split dd/di & m & median, IQR, 5th percentile; near-miss rate $[d_{\min} < r_{\mathrm{safe}}]$ & $\dt$, intruder aim and speed, prediction error; one bad run dominates a mean\\
  \midrule
  \multicolumn{5}{@{}l}{\emph{Efficiency}}\\
  Path length & $L_i = \sum_t \norm{\pos_i(t{+}1) - \pos_i(t)}$; ratio $\rho_i = L_i/L_i^{\mathrm{nom}}$ & m, -- & mean of $\rho_L$ with $t$ interval; paired difference to a baseline & controller jitter and noise inflate $L_i$; arena size (use the ratio)\\
  Travel time & $T_i$ of \cref{eq:ch25-travel} & s & mean with $t$ interval; censored at the cap $T_{\max}$ & goal tolerance $\eps_g$, intruders that hit a hovering drone, the cap\\
  Makespan & $\max_i T_i$ & s & mean or median; report with the success rate & the slowest drone only; cap; nominal delays of the plan\\
  Sum of costs & $\sum_i T_i$ & s & mean with $t$ interval & $k$ (normalise per drone when $k$ varies); cap\\
  \midrule
  \multicolumn{5}{@{}l}{\emph{Effort}}\\
  Replanning count & $n_{\mathrm{rp}}$ (Replanning entries); $n_{\mathrm{av}}$ steps in Avoiding & -- & mean with interval; distribution across runs & trigger thresholds, hysteresis, prediction noise (thrashing)\\
  Computation time & $c_i(t)$ per decision; $\bar c$, $c_{99}$, misses, $C_{\mathrm{tot}}$ & ms, s & mean and 99th percentile; cactus plot of $C_{\mathrm{tot}}$ & machine, language, load, logging; comparable only within one machine and run\\
  \midrule
  \multicolumn{5}{@{}l}{\emph{Constraints}}\\
  Formation error & $e_F(t)$ of \cref{eq:ch25-ef}; $\bar e_F$, $e_F^{\max}$ & m & mean with interval; peak distribution & template choice (fixed or nominal), formation graph $E_F$\\
  Communication violations & $n_{\mathrm{cv}}$ of \cref{eq:ch25-commviol} & steps & mean; fraction of steps; runs with $n_{\mathrm{cv}} > 0$ & $R_{\mathrm{comm}}$, kept edges, $k$; the nominal plan may already violate\\
  \bottomrule
  \end{tabularx}
\end{table}

\begin{pitfall}[Collision rate without minimum separation]
A collision rate of $0/20$ says that the true rate is below about $16\,\%$
(\cref{thm:ch25-three}) and nothing more. Two strategies with $0/20$ each
can differ enormously in how close they came: in the mini-study of
\cref{sec:ch25-study} the local-only and the replan-only strategies both
score $0/20$ collisions with two drones, but their median minimum
separations are $1.10$~m and $1.65$~m, and $40\,\%$ of the local-only runs
are near misses against $10\,\%$. Always report $d_{\min}$ (median, 5th
percentile, near-miss rate) next to the collision rate, and use the
segment-based minimum, not the sampled one.
\end{pitfall}

\begin{pitfall}[Comparing computation times across machines]
A decision that takes $0.17$~ms in the toy simulator of this chapter takes a
different time on a flight computer, under a different \python version,
with a different NumPy, or with the logger switched on. Computation times
are only comparable \emph{within} one study on one machine, and even then
they should be reported with the hardware and software
(\cref{sec:ch25-reporting}). When you compare with a published method,
compare counts that do not depend on the machine --- expanded nodes,
replans, low-level searches --- or run the published code yourself on your
machine. Never copy a runtime from a paper into your table.
\end{pitfall}

% ---------------------------------------------------------------------
\section{The experiment matrix and the baselines}
\label{sec:ch25-design}

\subsection{Factors and levels}

The training plan names five \textbf{factors}\index{experiment matrix} to
vary, each with a small number of \textbf{levels}
(\cref{tab:ch25-matrix}). A \textbf{cell} of the matrix is one choice of a
level for every factor; the mini-study of \cref{sec:ch25-study} occupies
four cells (two levels of the drone count, two of the intruder count, the
other three factors fixed). The \textbf{full factorial}\index{full factorial design}
design runs every cell. With the levels of \cref{tab:ch25-matrix} that is
$3 \times 4 \times 4 \times 3 \times 4 = 576$ cells, and with the $20$ seeds
per cell that \cref{sec:ch25-statistics} recommends as a minimum,
$11\,520$ runs. At a few seconds of simulated flight each this is a day of
computing for the toy simulator and weeks for a physics simulator --- and
most of the cells answer no question anyone asked.

\begin{table}[tb]
  \centering\small
  \caption{The experiment matrix of the training plan. The base cell (bold)
  is the reference around which a fractional design varies one factor at a
  time; the two starred factors interact and deserve a small full factorial
  of their own.}
  \label{tab:ch25-matrix}
  \begin{tabularx}{\textwidth}{@{}lY@{}}
  \toprule
  Factor & Levels\\
  \midrule
  Controlled drones $k$ & 2, \textbf{4}, 8 (more if computationally feasible)\\
  Unknown intruders $m$ & 0, \textbf{1}, 2, several crossing trajectories\\
  Prediction quality$^*$ & perfect state, \textbf{noisy state (KF)}, constant-velocity prediction, learned prediction (\cref{ch:ch20})\\
  Avoidance strategy$^*$ & local-only, replan-only, \textbf{hybrid} (plus the sanity baseline ``none'')\\
  Constraints & \textbf{no formation requirement} vs formation-preserving; \textbf{no communication radius} vs bounded range $R_{\mathrm{comm}}$\\
  \bottomrule
  \end{tabularx}
\end{table}

A \textbf{fractional design}\index{fractional design} runs a chosen subset
of the cells. The simplest and most readable one is
\emph{one-factor-at-a-time}: fix a base cell (bold in \cref{tab:ch25-matrix})
and vary one factor over all its levels while the others stay at the base.
With the levels above this costs $3 + 4 + 4 + 3 + 4 - 4 = 14$ cells, each of
them a clean answer to one question (``what happens as $k$ grows?''). Its
weakness is that it cannot see \emph{interactions}: whether the value of the
hybrid strategy depends on the prediction quality, for instance. When two
factors are suspected to interact, run their full factorial at the base
levels of the others --- strategy $\times$ prediction quality is
$3 \times 4 = 12$ cells --- and keep one-factor-at-a-time for the rest. State
the design in the paper; a reader who is told ``14 cells around a base
configuration plus the $3 \times 4$ strategy--prediction factorial'' can
judge what was and was not tested.

\subsection{Baselines}

Every number needs a reference, and the reference must be run on the same
scenarios (\cref{sec:ch25-scenarios}). \Cref{tab:ch25-baselines} lists the
five baselines\index{baseline (evaluation)} of this chapter. The three
avoidance strategies are the ones the plan names. ``No avoidance'' is a
\emph{sanity} baseline: it must collide often when an intruder is aimed at a
drone and never when there is no intruder, because the nominal plan is
conflict-free; if either fails, the scenario generator or the plan is broken,
and the whole study is meaningless. The fifth baseline is an
\emph{oracle}\index{oracle baseline}: on tiny instances, a mixed-integer
program (\cref{ch:ch22}) that knows the intruder's entire trajectory in
advance computes the shortest collision-free mission offline. No online
strategy can beat it, so it turns a path-length ratio into a statement about
the gap to the optimum instead of a comparison among imperfect methods. It is
expensive --- minutes per instance for two drones --- so it is run on a
handful of scenarios and reported as a separate line. The straight line of
the nominal plan is the trivial lower bound that costs nothing, and it is the
denominator of $\rho_L$.

\begin{table}[tb]
  \centering\small
  \caption{Baselines. Every one of them runs on exactly the scenarios of the
  strategy under test.}
  \label{tab:ch25-baselines}
  \begin{tabularx}{\textwidth}{@{}lYl@{}}
  \toprule
  Baseline & What it does & Purpose\\
  \midrule
  No avoidance & follows the nominal plan of \cbs/\ecbs and ignores intruders & sanity: collides with aimed intruders, never without\\
  Local-only & \orca/DWA with horizon $\ttc$, never replans; rejoins the moving reference & lower bound on effort, reference for deviation\\
  Replan-only & no reactive layer; \dstarlite/\astar replan whenever a conflict enters the horizon & reference for replanning count and path quality\\
  Hybrid & the state machine of \cref{ch:ch24}: local first, replan when reconnection fails & the method under test\\
  MILP oracle & offline optimum with full knowledge of the intruder (\cref{ch:ch22}), tiny instances & gap to the optimum on path length and time\\
  \bottomrule
  \end{tabularx}
\end{table}

\begin{pitfall}[Tuning on the test scenarios]
The hybrid strategy has thresholds --- the safety horizon $\tau_h$, the
deviation that triggers a replan, the minimum gap between replans --- and
so do the baselines. If you adjust them while watching the results table,
the table measures your patience, not the method. Split the seeds: tune
on seeds $0$--$9$, freeze every parameter in the configuration file, and
report seeds $100$--$119$ that no one has looked at. Tune the baselines
with the same care as your own method, on the same seeds, and say so; a
baseline that was ``run with default parameters'' is a straw man.
\end{pitfall}

% ---------------------------------------------------------------------
\section{Scenario generation}
\label{sec:ch25-scenarios}

\subsection{What a scenario contains}

\begin{definition}[Scenario]
\label{def:ch25-scenario}
A \textbf{scenario}\index{scenario} is the complete input of a run except
the strategy: the map, the starts $\vect{s}_i$ and goals $\vect{g}_i$ of the
$k$ drones, the nominal plan, the trajectories $\vect{q}_j(t)$ of the $m$
intruders for the whole horizon, the measurement noise sequence the
predictor will see, and the constraint data (formation template, kept edges,
$R_{\mathrm{comm}}$). It is a deterministic function
$\EvalScenario(\text{cell}, \text{seed})$ of the cell and an integer seed.
\end{definition}

The definition puts the noise into the scenario on purpose. If the predictor
draws its measurement noise from a global random state, two strategies see
different measurements and part of the difference between them is noise; if
the noise is a stored array indexed by time, both see the same. The same
holds for any randomness inside a strategy (the samples of \rrt, the
particles of a filter): give each drone its own generator seeded from the
scenario seed and the drone index. A scenario carries a
\textbf{signature}\index{scenario!signature}, a hash of its arrays, which
the runner uses to assert that every strategy of a cell and seed really
received the same object.

\subsection{Maps, starts and goals}

For the grid layer, take maps and start--goal sets from the MAPF benchmark
in the \code{.map} and \code{.scen} formats of \cref{ch:ch07}, so that the
global planner is tested on instances other people can download, and name
the map and the scenario file in the report. Random maps with an obstacle
density and random start--goal pairs are the alternative when you need a
knob for difficulty; then the map is part of the scenario and is regenerated
from the seed. The toy world of this chapter is obstacle-free: $k$ drones
start on the left edge of a $20 \times 20$~m arena, spread over $8$~m, and
fly to goals on the right edge whose rows are a random permutation of the
start rows, so that their straight-line missions cross. The nominal plan is
a prioritised plan (\cref{ch:ch08}) that gives each drone the smallest
departure delay keeping it $1.5$~m from every higher-priority drone; it
stands in for \cbs, whose output has the same shape (time-indexed,
conflict-free), and the code says so in its docstring.

\subsection{Intruder trajectory families and difficulty}

An intruder that wanders randomly rarely threatens anyone, and a study of
avoidance then measures nothing. The generator therefore \emph{aims}: it
picks a target drone $i$ and an instant $t_\times$ in the middle third of
that drone's mission, computes where the nominal plan puts the drone,
$\pos_\times = \pos_i^{\mathrm{nom}}(t_\times)$, and launches the intruder so
that it passes through $\pos_\times$ at $t_\times$ at a chosen
\textbf{crossing angle}\index{crossing angle} $\theta$ to the drone's heading
and a chosen speed. Three \textbf{trajectory families}\index{intruder!trajectory families}
are drawn around this aim (\cref{fig:ch25-intruders}):
\begin{itemize}
  \item \textbf{straight crossing}: constant velocity; a constant-velocity
        predictor is exact and the case measures the reaction itself;
  \item \textbf{turning}: constant turn rate $\omega$ ($5$--$15^\circ$/s),
        anchored so that the arc still passes the aim point; a
        constant-velocity prediction over $3$~s is off by
        $\tfrac12 v \omega \tau_h^2 \approx 1$~m;
  \item \textbf{evasive}: straight until $1$--$2.5$~s before $t_\times$, then
        a sudden heading change of $30$--$70^\circ$; the swerve is
        unpredictable, and it may take the intruder away from the target
        or into another drone.
\end{itemize}
Two knobs control difficulty. The \textbf{density} is the lateral spread of
the starts relative to $k$: with $8$ drones in $8$~m the nominal plan needs
long delays and the local layer sees several neighbours at once. The
crossing angle sets the geometry: $90^\circ$ is a broadside crossing,
$180^\circ$ head-on with the largest closing speed, $30^\circ$ an
overtaking intruder that stays in the drone's blind spot of a truncated
velocity obstacle. Report which angles and families a study used; ``random
intruders'' is not a description.

\begin{figure}[tb]
  \centering
  \inputfigure{ch25/intruders}
  \caption{Aimed intruders. The drone (blue) follows its nominal plan and
  reaches $\pos_\times$ at $t_\times$. All three families pass near
  $\pos_\times$ at $t_\times$: the straight crossing at angle $\theta$, the
  turning intruder along an arc anchored at the aim point, and the evasive
  intruder, which swerves shortly before $t_\times$ so that its
  constant-velocity prediction (dotted) points to the wrong place.}
  \label{fig:ch25-intruders}
\end{figure}

\subsection{The paired design}

\begin{definition}[Paired design]
\label{def:ch25-paired}
An evaluation is \textbf{paired}\index{paired design} if every strategy is
run on the same set of scenarios $\set{\EvalScenario(\text{cell}, s) :
s = 1, \dots, N}$, and every comparison between two strategies $A$ and $B$
is made through the per-scenario differences $D_s = X^A_s - X^B_s$ of a
metric $X$.
\end{definition}

\begin{proposition}[Pairing reduces variance when scenarios matter]
\label{thm:ch25-pairing}
Let $X^A_s$ and $X^B_s$ be the values of a metric for two strategies on
scenario $s$, with scenarios drawn independently, variances $\sigma_A^2$,
$\sigma_B^2$ and correlation $\rho$ between $X^A_s$ and $X^B_s$ across
scenarios. The paired estimate $\bar D = \frac1N \sum_s (X^A_s - X^B_s)$ of
the mean difference has variance
$(\sigma_A^2 + \sigma_B^2 - 2\rho\,\sigma_A\sigma_B)/N$, whereas the
unpaired estimate from two independent sets of $N$ scenarios has variance
$(\sigma_A^2 + \sigma_B^2)/N$. Pairing is more precise exactly when
$\rho > 0$, and for $\sigma_A = \sigma_B$ it shrinks the variance by the
factor $1 - \rho$.
\end{proposition}
\begin{proof}
For one scenario, $\Var(X^A_s - X^B_s) = \Var(X^A_s) + \Var(X^B_s) -
2\Cov(X^A_s, X^B_s) = \sigma_A^2 + \sigma_B^2 - 2\rho\sigma_A\sigma_B$, and
the mean of $N$ independent such differences divides the variance by $N$.
In the unpaired design the two sample means are independent, so their
variances add: $(\sigma_A^2 + \sigma_B^2)/N$. The difference between the two
expressions is $2\rho\sigma_A\sigma_B/N$, positive exactly when
$\rho > 0$.
\end{proof}

Scenarios matter: a head-on intruder at $2$~m/s lengthens every strategy's
path, an intruder that misses shortens none. So $\rho$ is positive, usually
strongly so, and the paired design sees a difference of a few per cent that
an unpaired design would bury under the scenario-to-scenario spread. In the
mini-study the path-length ratios of the hybrid and the local-only
strategies have $\rho = 0.86$ across the twenty scenarios of the four-drone
cell, so pairing cuts the variance of their comparison by a factor of
seven. \Cref{alg:ch25-paired} is the runner that enforces the design; it is
short because the work is in the scenario generator and the metric
functions.

\begin{algorithm}[htb]
\caption{Paired evaluation of a set of strategies on one cell}
\label{alg:ch25-paired}
\KwIn{cell $\kappa$ (levels of the five factors), seeds $s = 1, \dots, N$, strategies $\mathcal{S}$, configuration file $\mathcal{F}$}
\KwOut{one metric row per (seed, strategy), the raw logs, and per-strategy summaries}
\Fn{\EvalPairedStudy{$\kappa$, $N$, $\mathcal{S}$, $\mathcal{F}$}}{
  rows $\gets \emptyset$\;
  \For{$s \gets 1$ \KwTo $N$}{
    $\sigma \gets$ \EvalScenario{$\kappa$, $s$}\tcp*{deterministic in $(\kappa, s)$}\label{alg:ch25-paired:scenario}
    \ForEach{strategy $A \in \mathcal{S}$}{
      $\ell \gets$ \EvalRun{$\sigma$, $A$, $\mathcal{F}$}\tcp*{log: positions, states, $c_i(t)$, events}\label{alg:ch25-paired:run}
      save $\ell$ to disk under $(\kappa, s, A, \text{git hash of the code})$\;
      \textbf{assert} signature$(\sigma)$ unchanged\tcp*{the strategy must not mutate the scenario}\label{alg:ch25-paired:assert}
      rows $\gets$ rows $\cup \set{(\kappa, s, A, \EvalMetrics(\ell))}$\;
    }
  }
  \ForEach{metric $X$ and strategies $A, B \in \mathcal{S}$}{
    \EvalCompare{$(X^A_s)_s$, $(X^B_s)_s$}\tcp*{differences $D_s$, interval, effect size, tests}\label{alg:ch25-paired:compare}
  }
  \KwRet rows and summaries\;
}
\end{algorithm}

\paragraph{Walkthrough.}
Line~\ref{alg:ch25-paired:scenario} builds the scenario once per seed;
line~\ref{alg:ch25-paired:run} runs each strategy on it, and
line~\ref{alg:ch25-paired:assert} checks afterwards that the strategy did
not modify it (a replanner that edits the intruder array it was handed is a
bug that silently unpairs the design). The logs are saved before the
metrics are computed, so that a metric can be redefined later without
rerunning anything. Line~\ref{alg:ch25-paired:compare} does every comparison
on the differences $D_s$; the summaries per strategy (means, medians,
rates) are computed from the same rows. The same loop wrapped in an outer
loop over cells is the whole study; \code{run\_study()} in
\code{ch25\_evaluation.py} is this algorithm.

\begin{pitfall}[Unpaired comparisons]
The commonest flaw in student evaluations is to run strategy $A$ on twenty
random scenarios, then strategy $B$ on twenty \emph{other} random
scenarios, and compare the means. The scenario draw is then a hidden
factor: if $B$ happened to get an easier intruder in two of its runs, $B$
wins. \Cref{thm:ch25-pairing} says how much precision this throws away,
but the worse effect is bias in the small samples that are typical of
simulation studies. Seeds are free; use the same ones.
\end{pitfall}

% ---------------------------------------------------------------------
\section{Statistics: from runs to claims}
\label{sec:ch25-statistics}

\subsection{How many runs}

The number of runs $N$ per cell decides what the study can see. For a mean
with standard deviation $\sigma$ the half-width of a $95\,\%$ interval is
about $2\sigma/\sqrt{N}$: halving it costs four times the runs. For a rate
the situation is harsher, because rare events need many runs to occur at
all.

\begin{proposition}[The rule of three]
\label{thm:ch25-three}
If $N$ independent runs show no collision, the one-sided $95\,\%$ upper
confidence bound on the collision probability $p$ is
$1 - 0.05^{1/N} \approx 3/N$.
\end{proposition}
\begin{proof}
The probability of seeing no collision in $N$ runs is $(1-p)^N$. The
largest $p$ that is still consistent with this observation at level
$95\,\%$ satisfies $(1-p)^N = 0.05$, that is $p = 1 - 0.05^{1/N}$. Since
$0.05^{1/N} = e^{\ln(0.05)/N} \approx 1 + \ln(0.05)/N$ for large $N$ and
$-\ln 0.05 = 2.996$, the bound is approximately $3/N$.
\end{proof}

With $N = 20$ the bound is $0.139$ (the Wilson interval of the next
subsection gives $0.161$): a perfect score of $0/20$ rules out collision
rates above about one in six, and no more. To claim a collision rate below
$1\,\%$ you need at least $300$ clean runs; to \emph{measure} a rate of
$1\,\%$ with any precision, thousands. This is why safety is reported with
the minimum separation as well (\cref{sec:ch25-metrics}): $d_{\min}$ is
observed in every run and carries information even when nothing collides.
For the continuous metrics, $N = 20$ seeds per cell is a reasonable
minimum for a first study, $N = 50$--$100$ for a paper.

\subsection{Mean and standard deviation, or median and IQR}

The mean and the sample standard deviation
$s = \sqrt{\frac{1}{N-1}\sum_s (x_s - \bar x)^2}$ describe a symmetric,
light-tailed distribution well: path-length ratios, travel times without
timeouts, formation errors. The median and the interquartile range
describe skewed or heavy-tailed ones: minimum separations (bounded below by
zero, with a tail of near misses), computation times (a few slow decisions),
makespans with an occasional very late drone. The rule is simple: plot the
distribution first (a box plot, \cref{fig:ch25-boxplot}), and if the mean
and the median disagree visibly, report the median and say why. In the
mini-study the mean makespan of the local-only strategy with two drones is
$14.4$~s with a $t$ interval from $12.3$ to $16.6$~s, while the hybrid
strategy has $14.4$~s within $13.9$ to $15.0$~s: the means agree, but the
wide interval of the first hides one run in which a drone was pushed off
its route and needed $35$~s. The paired sign test of that comparison, which
sees the run-by-run order and not the magnitudes, says the hybrid is
\emph{faster} in one run and slower in ten.

\subsection{Confidence intervals}

Three intervals cover almost every need.\index{confidence interval}
\begin{itemize}
  \item The \textbf{$t$ interval} for a mean: $\bar x \pm t_{N-1,\,0.975}\,
        s/\sqrt{N}$, with $t_{19,0.975} = 2.093$ for $N = 20$. It assumes
        an approximately normal sampling distribution of the mean, which
        the central limit theorem grants for $N \ge 20$ unless the tail is
        extreme.
  \item The \textbf{bootstrap interval}\index{bootstrap} for any statistic
        \cite{efron1993bootstrap}: resample the $N$ values with replacement
        $B = 2000$ times, compute the statistic each time, and take the
        2.5th and 97.5th percentiles. It needs no distributional assumption
        and works for medians, percentiles and ratios; its randomness is
        seeded like everything else.
  \item The \textbf{Wilson interval}\index{Wilson interval} for a
        proportion $\hat p = x/N$ \cite{wilson1927probable}:
        \begin{equation}
          \frac{\hat p + \frac{z^2}{2N} \pm z\sqrt{\frac{\hat p(1-\hat p)}{N}
          + \frac{z^2}{4N^2}}}{1 + \frac{z^2}{N}}, \qquad z = 1.96.
          \label{eq:ch25-wilson}
        \end{equation}
        Unlike the textbook $\hat p \pm z\sqrt{\hat p(1-\hat p)/N}$, it does
        not collapse to a width of zero at $\hat p = 0$: for $0/20$ it gives
        $[0, 0.161]$, for $1/20$ $[0.009, 0.236]$, for $17/20$
        $[0.640, 0.948]$.
\end{itemize}

\subsection{Paired comparisons and effect sizes}

Given the per-scenario differences $D_s = X^A_s - X^B_s$, report three
things.\index{paired comparison}
\begin{enumerate}
  \item \textbf{The mean difference with its $t$ interval}, in the units of
        the metric. ``The hybrid strategy's paths are $2.1\,\%$ longer than
        local-only, $95\,\%$ interval $0.6$ to $3.6\,\%$'' is a complete
        statement: it gives the size, the direction and the uncertainty. If
        the interval excludes zero, the difference is significant at the
        $5\,\%$ level; but the interval says more than a $p$-value does.
  \item \textbf{A standardised effect size}\index{effect size}: Cohen's
        $d_z = \bar D / s_D$, the mean difference in units of the standard
        deviation of the differences \cite{cohen1988power}. It lets you
        compare metrics with different units and studies with different
        $N$; values around $0.2$, $0.5$ and $0.8$ are conventionally called
        small, medium and large. Report it next to the raw difference, not
        instead of it.
  \item \textbf{A rank test} when the differences are skewed or $N$ is
        small: the Wilcoxon signed-rank test \cite{wilcoxon1945individual},
        which ranks $\abs{D_s}$ and asks whether the positive ranks
        dominate, and the sign test, which only counts how often $A$ beat
        $B$. The counts ``better / worse / tie'' are worth printing in any
        case; a reader understands ``16 of 20 scenarios'' faster than any
        $p$-value.
\end{enumerate}
Every comparison you run is a chance to be fooled by noise. With $10$
metrics and $3$ strategy pairs you run $30$ tests, and at the $5\,\%$ level
one or two of them will be ``significant'' by accident. Either decide in
advance which two or three comparisons are the primary ones and treat the
rest as exploratory, or adjust: Holm's procedure \cite{holm1979simple}
sorts the $M$ $p$-values and compares the $j$-th smallest with $0.05/(M-j+1)$,
a step-down version of the Bonferroni correction that loses little power.
\Cref{tab:ch25-stats} maps the questions of a results section to these
tools.

\begin{table}[tb]
  \centering\small
  \caption{Which statistic answers which question. All comparisons are
  paired.}
  \label{tab:ch25-stats}
  \begin{tabularx}{\textwidth}{@{}YlY@{}}
  \toprule
  Question & Statistic & Plot\\
  \midrule
  How often does it collide? & rate with Wilson interval; rule of three for $0/N$ & bar with interval; Pareto plot\\
  How close does it come? & median and 5th percentile of $d_{\min}$; near-miss rate & box plot per level\\
  How much longer are the paths? & mean $\rho_L$ with $t$ interval; paired difference & box plot; Pareto plot\\
  Is $A$ better than $B$? & mean of $D_s$ with interval, $d_z$, better/worse/tie, Wilcoxon or sign test & paired dots or a difference box plot\\
  How does it scale with $k$? & one-factor-at-a-time cells; success rate under a cap & line per strategy against $k$; cactus plot\\
  Does it meet its deadline? & $c_{99}$, deadline misses & cactus/survival plot of $C_{\mathrm{tot}}$ or $c_i(t)$\\
  Which trade-off is best? & pairs (safety, efficiency) per cell and strategy & Pareto plot with the front marked\\
  \bottomrule
  \end{tabularx}
\end{table}

\subsection{Success rate with a runtime cap}

Some runs do not finish: a planner exceeds its time budget, a drone never
reaches its goal, the simulation hits the mission cap $T_{\max}$. The MAPF
literature reports the \textbf{success rate}\index{success rate} --- the
fraction of instances solved within the cap --- as a function of the number
of agents \cite{sharon2015cbs,barer2014ecbs,stern2019mapf}, and every
metric that needs a finished run (path length, makespan) is averaged over
the successful runs only, \emph{with the success rate printed next to it}.
Two rules keep this honest. First, when two strategies are compared on a
finished-run metric, compare them on the scenarios that \emph{both}
finished, or the better strategy is penalised for finishing the hard cases.
Second, when a single number is needed, use a penalised average such as
PAR-2\index{PAR-2 score}: an unfinished run counts as twice the cap. The
cactus plot of \cref{fig:ch25-cactus} avoids the problem altogether by
showing the whole distribution of the effort.

\begin{pitfall}[Reporting means without spread, and hiding failed runs]
A table of means with no interval, no $N$ and no success rate is not a
result; it is an anecdote with decimals. Every entry needs its interval or
IQR and the number of runs it is based on, and every cell needs its success
rate. Runs that crashed, timed out or collided are the most informative
ones: keep them in the raw results file, count them in the table, and
describe the worst of them in the text. Dropping a ``pathological seed''
after the fact is the same thing as tuning on the test set.
\end{pitfall}

\subsection{Three plots}

\paragraph{Box plots per factor level.}
A \textbf{box plot}\index{box plot} shows the median (line), the
interquartile range (box) and the whiskers at $1.5$ IQR, with outliers as
dots: five numbers that a mean and a standard deviation cannot replace for
a skewed metric. Draw one box per strategy for every level of the factor
under study, as in \cref{fig:ch25-boxplot}, and put the safety threshold
($r_{\mathrm{safe}}$, the collision distance) as a horizontal line so that
the eye can count the boxes that cross it. In \cref{fig:ch25-boxplot} the
local-only and hybrid boxes sit on the safety radius $r_{\mathrm{safe}} =
1$~m, because a velocity-obstacle avoider chooses velocities on the boundary
of the forbidden set, while the replan-only boxes sit $0.6$~m higher, because
its detours are offset by at least $2$~m; the path-length panel shows the
price.

\begin{figure}[tb]
  \centering
  \inputfigure{ch25/boxplot}
  \caption{Box plots from the mini-study of \cref{sec:ch25-study} (one
  intruder, $20$ seeds per box). Left: minimum drone--intruder separation
  $d_{\min}^{\mathrm{di}}$; the dashed line is the safety radius
  $r_{\mathrm{safe}} = 1$~m and the dotted line the collision distance
  $0.6$~m. Without avoidance every median is at zero. Right: path-length
  ratio $\rho_L$. The replan-only strategy buys its larger separation with
  the longest paths; the hybrid sits between the two.}
  \label{fig:ch25-boxplot}
\end{figure}

\paragraph{Cactus and survival plots for runtime.}
A \textbf{cactus plot}\index{cactus plot} sorts the runs of each strategy by
their effort (total computation time, or the number of expansions) and
plots, against the effort $x$, the number of runs that finished with effort
at most $x$. A curve further to the right and lower is a slower method; a
curve that stops before $N$ has unsolved runs, and the height at which it
stops is the success rate. Its mirror image, the \textbf{survival
plot}\index{survival plot}, shows the fraction of runs still unfinished at
effort $x$. Both show the whole distribution, including the tail that a mean
hides. \Cref{fig:ch25-cactus} draws the total computation time per run of
the four strategies across all $80$ one-intruder and no-intruder runs; the
no-avoidance curve is the cost of prediction and tracking alone, and the two
strategies with a reactive layer pay for their velocity sampling in every
Avoiding step, whereas replanning is cheap here because a replan happens a
few times per run.

\begin{figure}[tb]
  \centering
  \inputfigure{ch25/cactus}
  \caption{Cactus plot of the total computation time per run
  ($C_{\mathrm{tot}}$, milliseconds of \python on one machine) for the
  $80$ runs of each strategy in the mini-study. Every curve reaches $80$:
  no run exceeded the cap. The reactive strategies are the slowest here
  because their velocity sampling runs at every Avoiding step; the
  absolute values mean nothing outside this machine, the ordering does.}
  \label{fig:ch25-cactus}
\end{figure}

\paragraph{Pareto plots of safety against efficiency.}
Safety and efficiency pull in opposite directions, and a strategy that is
best in one column and worst in another is not ``better''. A
\textbf{Pareto plot}\index{Pareto plot} puts one point per strategy and
cell into the plane (efficiency loss, safety loss) --- here $\rho_L$ against
the near-miss rate or the collision rate --- and marks the points that are
not dominated by any other. In \cref{fig:ch25-pareto} the three avoidance
strategies form a front: local-only is the most efficient, replan-only the
safest, the hybrid in between, and no-avoidance is dominated by all of them.
The plot states the trade-off; the text must then say which point the
application wants.

\begin{figure}[tb]
  \centering
  \inputfigure{ch25/pareto}
  \caption{Pareto plot of the mini-study: mean path-length ratio (efficiency
  loss, right is worse) against the near-miss rate (safety loss, up is
  worse), one marker per strategy and cell; filled markers are the
  two-drone cells, hollow markers the four-drone cells, labels give the
  collision rate. No-avoidance (top left) is dominated; the three avoidance
  strategies form the front.}
  \label{fig:ch25-pareto}
\end{figure}
